\section{Relationship to existing evidence and methods}
\label{sec:related}

The returned-gradient case and gradient-gain study supply primary evidence
about observation coverage and scoring units \citep{paper1,gradients2026}.
The latter's exact-document recovery rises from $13.71\%$ to $37.77\%$
in its controlled GPT-2 comparison; that does not mean a comparable
percentage-point gain in token recovery. These published findings are
reused as source evidence, not presented as new independent experiments.

The LLM-Attack Suite already specifies observation contracts, controls,
saved predictions, missing-attack disclosures, and measurement-conformance
checks \citep{suite2026}. Our contribution is a retrospective reconciliation
of claims across campaigns and abstraction levels. The contract terminology
is a reporting device; the five outcomes are not claimed as a novel
security logic with independently measured diagnostic accuracy. This paper
does not compare a new attack suite against the existing one or claim
superior attack coverage.

\citet{privpair2026} studies split federated LLM fine-tuning under a
server observing activations and gradients, argues that perturbation-only
protection is insufficient in its setting, and proposes client-side
obfuscate-and-recover adapters. That is close training-side context,
rather than evidence for a universal value/structure dichotomy.
\citet{mia2026} separately demonstrates evaluator and attack dependence
in image reconstruction. Its findings motivate careful evaluation; they
do not establish that qualified breaches in R1 or Paper 4 are instrument
artifacts. MI-estimator limitations \citep{mcallester2020,song2020}
constrain critic-score interpretation irrespective of the defense.

\section{Limitations and implications for further evaluation}
\label{sec:limits}

\paragraph{One program and retrospective selection.}
The cases were selected from the evidence underlying the original
synthesis, after their outcomes were known. Shared code, corpora, and
research decisions create dependence. Case counts do not estimate
field-wide failure prevalence; correcting known examples does not measure
the sensitivity or specificity of the audit on unseen studies. A
prospectively specified independent cohort would be needed for that claim.

\paragraph{Scope of recomputation.}
We reconstruct selected statistics and decisions from supplied artifacts,
not the original GPU execution. F1 and parts of F2 remain source-reported
through published or prior audit records; the new script does not recover
every frame or retrain their attackers. Source hashes identify consumed
files, not the truth of all uninspected collection history. Records lacking
an admissible primary derivation are excluded from quantitative headlines
rather than filled with plausible values.

\paragraph{Calibration, uncertainty, and power.}
Qualification applies to a specified binding, not an unbounded adversary.
The four-seed R1 and three-seed pilot summaries are conditional on fixed
evaluation panels and have limited precision. There is no newly executed
confirmatory study. The separate power calculation is conditional planning,
not evidence of universal impossibility. Policy thresholds, positive-control
floors, and the smallest reliably detectable injected leak need different
definitions and units.

\paragraph{Security and utility.}
Non-detection is finite and attack-relative. A low pooled recovery score,
an invertible shared-key wrapper, or a near-zero lower-bound estimate cannot
stand in for compromised-host privacy. Restoration on tested prompts does
not establish universal zero utility cost, and short-context NLL does not
capture every quality endpoint. Noise and quantization require specified
strengths and utility budgets; the records do not justify a theorem that
all value perturbation is futile.

The practical implication is to retain the experimental result while
controlling the claim it supports. An independent validation could test
whether these distinctions identify known evaluation faults without
dismissing correctly detected breaches. Its data, budgets, attack
adaptation, and primary contrast would need freezing before the validation
outcomes were observed. That proposed study is not counted as a result here.

\section{Conclusion}

The program contains an omitted-channel failure, an estimand mismatch,
incomplete instrument qualification, and valid negative defense results.
The latter should survive an audit. Reconstructing the decisions shows why
20 unqualified R1 instruments do not explain away 13 qualified breaches,
and why a seven-cell pilot with no eligible candidate does not imply a
broken gate. Scope checks prevent two further promotions: wire-observer
non-detection into compromised-host protection, and a near-zero MI lower
bound into a small leakage certificate. The reproducible contribution is
the separation of these admissible conclusions, with their sources and
limitations, rather than a universal defense or information-theoretic claim.

\paragraph{Data and code availability.}
The accompanying revision directory contains secondary-analysis code,
derived tables, vector-figure scripts, source-hash bindings, and a build
recipe. Primary inputs remain in the Paper 9 artifact collection under
the relative paths in Appendix~\ref{app:artifacts}. This text does not
assert a public repository URL or access to all original execution
environments. No new training or attack campaign is required to rerun the
packaged arithmetic; rerunning the historical GPU studies is a separate
reproduction task \citep{artifacts2026}.
